\documentclass[11pt]{article}
\usepackage[a4paper,margin=27mm]{geometry}
\usepackage{amsmath,amssymb,amsthm,hyperref}
\newtheorem{theorem}{Theorem}
\title{A sharp obstruction at the weak triple-minimum support bound}
\author{Research note: weighted Gram matrices and row leverage}
\date{30 September 2026}
\begin{document}
\maketitle

Let a probability distribution on $q$ permutations of $n$ labels have
positive row weights $w_1,\ldots,w_q$. Suppose each member of every
triple is its earliest label with probability $1/3$. Pairwise
unbiasedness is not assumed.

\begin{theorem}\label{support}
For $n\ge5$, such a distribution has $q\ge n$.
For $n=4$, the sharp bound is $q\ge3$.
\end{theorem}
\begin{proof}
For distinct labels $i,j$, put $p_{ij}=\Pr(i<j)$. The triple-minimum
condition implies $1/3\le p_{ij}\le2/3$.
Call $i$ bad if $p_{ij}=1/3$ for some opponent $j$, and good otherwise.
For a bad label, every event $i<j,i<k$ already has the same probability
as $i<j$. Thus $i<j$ forces $i<k$ almost surely for every other label
$k$, and the bad label is globally first with probability exactly
$1/3$. There are at most three bad labels.

Fix a pivot $i$ and let $C$ be the $q\times(n-1)$ zero--one matrix
whose opponent column $j$ records the event $i<j$. Write
$W=\operatorname{diag}(w_1,\ldots,w_q)$, $e_j=p_{ij}-1/3$, and
$D=\operatorname{diag}(e_j)$. Then
\begin{equation}\label{gram}
 G=C^{\mathsf T}WC=D+\tfrac13\mathbf1\mathbf1^{\mathsf T},
 \qquad C^{\mathsf T}w=e+\tfrac13\mathbf1.
\end{equation}
If $q<n-1$, all $e_j>0$ would make $G$ have rank $n-1>q$.
Hence every pivot would be bad, forcing $n\le3$. This proves
$q\ge n-1$ when $n\ge4$.

It remains to exclude $q=n-1$ for $n\ge5$. For a good pivot all
$e_j>0$, so $G$ and the now square matrix $C$ are invertible.
The solution $v=C^{-1}\mathbf1$ can be found from \eqref{gram}:
Sherman--Morrison gives
\begin{equation}\label{v}
 v_j=1-\frac{\alpha}{e_j},\qquad
 \alpha=\frac{(q-1)/3}{1+\frac13\sum_j e_j^{-1}}>0.
\end{equation}
Thus $v_j<1$ for every $j$. A row of $C$ with zero or one entries
equal to one cannot have scalar product one with $v$. Every good
pivot therefore has at least two labels after it in every row.

The last two positions of every permutation must consequently be
bad labels. If there were at most two bad labels, they would occupy
these positions in every row and never be first, a contradiction.
Thus there are exactly three bad labels, each first with probability
$1/3$, and $q-2$ good labels.

For a fixed good label $a$, sum its minimum probabilities over the
three pairs of bad labels. The total is one. In every row at least
one bad label precedes $a$, so a row contributes at most one to this
sum. Since all weights are positive, every row contributes one:
exactly one bad label precedes $a$. Hence every permutation consists
of its first bad label, then all good labels, then the other two bad
labels. In particular, for good $a$ and bad $b$,
\begin{equation}\label{badpair}
 p_{ab}=2/3.
\end{equation}

For two good labels $a,c$, let $r=p_{ac}$, and let $x_b$ be the
probability of $a<c$ in the rows beginning with bad label $b$.
Outside these rows, both $a,c$ precede $b$, so the minimum condition
on $\{a,c,b\}$ gives $r-x_b=1/3$. Summing over the three bad labels
gives $r=3(r-1/3)$, and therefore
\begin{equation}\label{goodpair}
 p_{ac}=1/2.
\end{equation}

Return to a good pivot's matrix. Equations \eqref{badpair} and
\eqref{goodpair} mean that $D$ has three entries $1/3$ (bad
opponents) and $q-3$ entries $1/6$ (good opponents). Put
$u=D^{-1}\mathbf1$, so these entries of $u$ are respectively three
and six. Direct inversion of \eqref{gram} gives
\[
 G^{-1}=D^{-1}-\frac{uu^{\mathsf T}}{6(q-1)},\qquad
 CG^{-1}C^{\mathsf T}=W^{-1}.
\]
Fix a row of weight $w_s$. If this good pivot has $\ell$ other good
labels after it, it also has exactly two bad labels after it.
The diagonal entry of the last matrix identity is thus
\begin{equation}\label{leverage}
 \frac1{w_s}=6t-\frac{6t^2}{q-1},\qquad t=\ell+1.
\end{equation}
In the same row, as the pivot ranges over all good labels, $t$ takes
every value $1,\ldots,q-2$. If $q\ge5$, comparison of $t=1$ and
$t=2$ in \eqref{leverage} would force $q=4$, a contradiction.
If $q=4$, \eqref{leverage} forces every row weight to be $1/4$.
No event can then have probability $1/3$, again a contradiction.
This excludes equality for every $n\ge5$.

For $n=4$, the three equally weighted permutations
$1203,3201,0213$ satisfy the required triple-minimum condition,
so the bound $q\ge3$ is sharp.
\end{proof}

The proof applies to arbitrary positive row weights. In the
equal-weight case one additionally has $3\mid q$, so
$q\ge3\lceil n/3\rceil$ for $n\ge5$.

The equality argument also describes every four-label, three-row
example: all weights are $1/3$; one label always occupies the second
position; each of the other three labels is first once, and the
remaining two positions can be ordered arbitrarily.

\paragraph{Verification and companion arguments.}
The weighted argument was independently derived and checked in two
parallel research threads. It uses no enumeration or optimization
solver. The earlier uniform-weight determinant proof is preserved
in \path{weak_minimum_uniform_determinant.tex}; it forces
$3h-1$ to be a square modulo three when $q=3h=n-1$, $h\ge2$.
The note \path{six_weak_minimum_permutations.tex} gives another
independent argument for $q=6,n=7$. Small exact profile calculations
are retained in \path{check_weak_minimum_profiles.py} and its JSON
output, but are not needed for Theorem~\ref{support}.
No claim of literature novelty is made here.
\end{document}
